\chapter*{Abstract}

In recent years, Large Language Models (LLMs) have transformed natural language processing (NLP), excelling in tasks like chatbots, summarization, and content analysis. However, they often struggle with domain-specific questions, producing inaccurate or irrelevant results, especially when data is outdated or sparse. This is particularly problematic in political discourse, where positions evolve over time.

Party manifestos are used to assess party ideologies, policy stances, and ideas. While LLMs work well with context, manifestos are lengthy and complex, and being in Portuguese, a language with fewer resources, may lead to even worse results.

To address these challenges, the Retrieval-Augmented Generation (RAG) framework was developed, integrating external data retrieval with LLMs to reduce unreliable information and enhance their practical applications.

This research seeks to evaluate whether a RAG-based system, \sysname{}, can improve the accessibility and comprehension of Portuguese political party programs, allowing citizens and researchers to better understand policy positions. By exploring RAG implementations, this study will explore the effectiveness, scalability, and practical application of these AI-driven methods for analyzing political texts.

\paragraph{Keywords} Natural Language Processing, Information Extraction, Large Language Models, Retrieval-Augmented Generation, Political texts

\cleardoublepage

\chapter*{Resumo}

Nos últimos anos, Large Language Models (LLMs) revolucionaram o processamento de linguagem natural (NLP), destacando-se em tarefas como chatbots, sumarização e análise de conteúdo. No entanto, estes modelos enfrentam frequentemente dificuldades com perguntas mais específicas, produzindo resultados imprecisos ou irrelevantes, especialmente quando os dados estão desatualizados ou são escassos. Este problema é particularmente relevante no discurso político, onde as posições evoluem ao longo do tempo.

Os manifestos partidários são utilizados para avaliar as ideologias, posições políticas e propostas dos partidos. Embora os LLMs funcionem bem com contexto, os manifestos são documentos extensos e complexos e, por estarem em português, uma língua com menos recursos, podem levar a resultados ainda menos satisfatórios.

Para enfrentar estes desafios, foi desenvolvido o framework Retrieval-Augmented Generation (RAG), que integra a recolha de dados externos com os LLMs, reduzindo a produção de informação pouco fiável e otimizando as suas utilizações práticas.

Esta investigação procura avaliar se um sistema baseado em RAG, o \sysname{}, pode melhorar a acessibilidade e a compreensão dos programas dos partidos políticos portugueses, permitindo que cidadãos e investigadores compreendam melhor as posições políticas. Ao explorar implementações de RAG, este estudo analisará a eficácia, escalabilidade e aplicação prática destes métodos impulsionados pela inteligência artificial para a análise de textos políticos.

\paragraph{Palavras-chave} Processamento de Linguagem Natural, Extração de Informação, aqui, separadas, por, vírgulas

\cleardoublepage
